% !TeX root = ../main.tex

\chapter{Conclusion and Future Work}
\label{ch:conclusion}

\section{Summary of Contributions}

This thesis systematically evaluates the suitability of general-purpose genomic foundation models (GFMs) for metagenomic short-read taxonomic classification. Through controlled ablations of data scale, encoder depth, pre-training, and tokenization, we identify which design factors are decisive and which are not. Our main findings are:

\textbf{Tokenization sets the ceiling; data volume dominates only inside a fixed tokenizer.} Through systematic data scaling spanning 500K to 50M training reads, we show that within the NT-v2 6-mer framework, increasing training data by $10\times$ ($500\text{K}\to 5\text{M}$) yields $+7.76$~pp in accuracy and $+12.27$~pp in macro F1; a further $10\times$ ($5\text{M}\to 50\text{M}$) adds $+4.02$~pp. Pushing a further $5\times$ to 250M reads adds only $+0.22$~pp (67.29\%), empirically refuting the log-linear projection: the NT-v2 6-mer model saturates beyond 50M. This saturation is tokenizer-specific rather than a data limit: trained on the same reads over the same $50\text{M}\to250\text{M}$ range, an overlapping 13-mer model scales from 87.5\% to 98.7\% genus Top-1. By contrast, architectural and training-technique ablations on 500K reads yield at most $\pm0.5$~pp. The token-level pipeline---attention-pooling head fine-tuned with LoRA on the NT-v2 backbone---reaches 67.07\% genus Top-1 RC TTA at 50M reads, compared to 51.70\% for the frozen mean-pool baseline.

\textbf{The $+13.19$~pp one-layer gap is depth, not pre-training; pre-training buys data efficiency.} A 29-layer pre-trained NT-v2 versus a one-layer randomly initialized Transformer at the same 50M scale and 6-mer tokenizer differs by $+13.19$~pp RC TTA (67.07\% vs.\ 53.88\%). That comparison confounds pre-training with depth and width. A depth sweep with the tokenizer, data, head, and \texttt{clean\_common} evaluation held fixed scores 53.92\% / 65.44\% / 67.94\% at 1 / 8 / 16 from-scratch layers, against 67.08\% for NT-v2 LoRA: depth accounts for $+14.02$~pp ($1\to 16$), and at 50M the net contribution of pre-training is $\approx 0$ (the 3.88M-parameter 16-layer model is $+0.86$~pp higher). The 8- and 16-layer models are trained 100\% from scratch, so the $\sim$68\% 6-mer ceiling is not an artefact of LoRA updating only $\sim$1\% of NT-v2. The ranking reverses at 5M reads under the same 16-layer match: NT-v2 63.05\% versus from-scratch 54.11\% ($+8.94$~pp). Pre-training therefore helps when labels are scarce; it does not raise the 6-mer ceiling once 50M labelled reads are available. Within the family of pre-trained GFMs at 5M reads, NT-v2 still outperforms DNABERT by $+1.27$~pp and DNABERT-2 by $+4.17$~pp.

\textbf{Tokenization remains the dominant design factor, including at matched architecture.} Holding MetaTransformer fixed and varying $k \in \{6, 12, 13\}$ and stride, overlapping stride$=$1 is critical (up to $+60$~pp over non-overlapping at $k=13$), and longer $k$ yields monotonically higher accuracy among overlapping configurations. MetaTransformer with overlapping 13-mer tokenization, trained from scratch on the same 50M reads, achieves \textbf{87.42\%} genus Top-1---\textbf{+20.35~pp} above NT-v2 + LoRA (67.07\%). A same-architecture 16-layer comparison on \texttt{clean\_common} RC TTA reaches 83.83\% with a hashed 13-mer (2~GiB embedding) and 85.63\% with an exact 13-mer, versus 67.94\% for the 16-layer non-overlapping 6-mer; a 29-layer from-scratch 6-mer reaches only 69.52\%. Even when the 6-mer model is given 16 layers and the 13-mer MetaTransformer only one, the 13-mer still wins by $\sim$20~pp (87.42\% vs.\ 67.94\%). Depth compensates for a weak tokenizer; it does not break the 6-mer ceiling. MetaTransformer's published advantage on this task is therefore its 13-mer tokenizer, not a uniquely strong encoder.

\textbf{RC TTA is a robust, free improvement.} Reverse-complement test-time augmentation consistently improves accuracy by $+0.1$ to $+1.5$~pp across experimental conditions, requiring no additional training cost. The benefit is largest for pre-trained backbones ($+0.78$ to $+1.54$~pp) and near-zero for randomly initialized models ($+0.08$~pp), supporting the interpretation that RC TTA primarily corrects pre-training-induced strand biases rather than stochastic inference noise.

\textbf{Class balance is critical for rare taxa.} Balanced subsampling across species reduces the number of F1$=$0 classes from 9 to 2 and improves macro F1 by $+12.27$~pp, demonstrating that class imbalance---not model capacity---is the primary cause of poor performance on rare taxa at the 500K--5M scales.

\textbf{Per-genus specialisation validates the hierarchical architecture; subgenus routing does not.} Evaluation on a 100K independent test set demonstrates that per-genus 50M LoRA classifiers under oracle-genus routing reach 29.5\% species Top-1, exceeding the flat sp\_v4 oracle ceiling of 27.4\%. Predicted routing through the 66.1\%--67.07\% genus model lowers Top-1 to 15.6\%, marginally below the flat baseline (15.9\%), because per-genus specialists do not recognise reads from genera they were not trained on. Subgenus $K$-means routing on NT-v2 embeddings is a negative result (oracle Top-1 25.6--26.1\%).

\textbf{Per-read accuracy translates to sample-level utility under controlled conditions, but genus abundance cannot choose a method.} On synthetic species-balanced random-partition samples, 67\% per-read accuracy yields Pearson $r = 0.993$ for relative abundance and Bray-Curtis dissimilarity 0.094; 100\% detection sensitivity is achieved for taxa at $\geq$1\% relative abundance. All 100 random-partition samples share the same expected per-genus abundance profile, so this result measures stability under binomial sampling noise rather than generalisation to genuinely variable communities. Kraken2+Bracken matches neural genus abundances ($r = 0.997$) on these simulated communities. Evaluation on real and compositionally diverse microbiome samples remains an open question.

\section{Limitations}

This work has several limitations that should be acknowledged:

\begin{enumerate}
    \item \textbf{Simulated data only}: All experiments use reads simulated by ART from reference genomes. Real metagenomic reads contain additional noise (host contamination, chimeric reads, strain-level variation) that may degrade classifier performance.
    
    \item \textbf{Limited taxonomic scope}: The reference database covers 120 genera and 1,535 species of human gut microorganisms. Performance on environmental samples with broader taxonomic diversity is unknown.
    
    \item \textbf{Single read length}: All experiments use 150~bp reads. Performance at other read lengths (e.g., 250~bp for Illumina MiSeq, or long reads from Oxford Nanopore / PacBio) has not been evaluated.
    
    \item \textbf{Data scaling saturates for the 6-mer model, not universally}: We trained balanced subsets up to 250M reads (Section~\ref{sec:exp-scaling}); the warm-started 250M NT-v2 model reaches only 67.29\% RC TTA ($+0.22$~pp over the 50M 67.07\%), confirming that the 6-mer model saturates rather than following the log-linear projection. This is a limit of the tokenizer, not the data: trained on the same reads over the same range, an overlapping 13-mer model scales from 87.5\% to 98.7\% genus Top-1. The gap to overlapping-$k$-mer methods is therefore not closable by additional training data within the NT-v2 6-mer framework. We note that the 13-mer model's 98.7\% reflects near-complete closed-set coverage of the 1{,}535 reference genomes' 13-mer content (approaching a lookup over known genomes), and does not measure generalisation to novel genomes.

    \item \textbf{Memory-bound loader in the primary pipeline}: The single-GPU TWCC pipeline loads training reads into memory; the 250M-scale experiments above were instead run on a multi-GPU H200 cluster with a lazy, index-based loader. Unifying the two (a streaming loader in the main codebase) remains future engineering work.
    
    \item \textbf{Architectural capacity ceiling of MetaTransformer}: MetaTransformer's single-encoder-block design (${\sim}$5M encoder parameters, excluding the embedding table) saturates within one training epoch on 50M reads; a second epoch produces overfitting rather than further generalisation. A larger model would likely improve beyond 78.83\% at 12-mer but was not evaluated due to infrastructure constraints.

    \item \textbf{Scope is general-purpose GFMs}: We evaluated NT-v2 and, at 5M scale, DNABERT and DNABERT-2---models pre-trained for eukaryotic or mixed-genome regulatory genomics, not for short-read metagenomic taxonomy. Two GFMs pre-trained on metagenomic sequence were not tested: METAGENE-1~\cite{Liu2025metagene} ($\sim$1.5~Tbp wastewater metagenomes, BPE) and GenomeOcean~\cite{Zhou2025genomeocean} ($\sim$600~Gbp metagenomic assemblies, 4B parameters, BPE). The domain and context-length mismatch that plausibly explains NT-v2's $\approx 0$ transfer at 50M would not apply to those models in the same way, and their BPE tokenizers would introduce a further variable. Claims in this thesis are therefore scoped to general-purpose GFMs, not to all GFMs.

    \item \textbf{Remaining confounds in the tokenizer comparison}: The same-architecture 16-layer 13-mer arms differ in exact versus hashed vocabulary, embedding width ($d_{\text{model}}$ 64 vs.\ 128), and canonical versus non-canonical vocabulary; Job~A ran fp16 and Jobs~B/C ran bf16. MetaTransformer's 87.42\% is quoted from the matched 50M protocol used in Chapter~\ref{ch:experiments}; a one-layer hashed 13-mer in our own codebase was not run. Soil depth-matched from-scratch controls were not completed.

    \item \textbf{LoRA versus full fine-tune}: NT-v2 is adapted with LoRA on $\sim$1\% of parameters. A full fine-tune might score higher. The from-scratch 8- and 16-layer models, which train 100\% of parameters and hit the same $\sim$68\% 6-mer ceiling, are the evidence that the ceiling is not solely a LoRA artefact.
\end{enumerate}

\section{Future Work}

The findings above point to four primary future directions:

\subsection{Microbial Foundation Models with Overlapping Long-$k$ Tokenization}

The tokenization ablation (Section~\ref{sec:exp-tokenization}) identifies overlapping long-$k$ tokenization as the dominant performance driver. No publicly available GFM that we evaluated combines this property with large-scale pre-training on diverse microbial genomes; METAGENE-1 and GenomeOcean use BPE rather than long $k$-mers and were not tested here. The most promising direction consistent with the measurements in this thesis is therefore to pre-train a transformer on the GTDB database~\cite{Parks2022} ($>$400,000 microbial genomes) with overlapping 12-mer or 13-mer tokenization---or to evaluate hashed long-$k$ embeddings, which reached 83.83\% at 16 layers with a 2~GiB table. Open empirical questions include (i)~whether RC TTA provides an equivalent benefit ($\sim+1$~pp) at $k = 12$/13, (ii)~whether $k > 13$ continues to improve accuracy, (iii)~whether a one-layer hashed 13-mer in our codebase matches MetaTransformer's 87.42\%, and (iv)~whether BPE models pre-trained on metagenomes (METAGENE-1, GenomeOcean) close the 50M gap that NT-v2 does not. A complementary, lower-cost direction is to fine-tune NT-v2 with a learned 12-mer projection layer; this is bounded above by the model's pre-trained 6-mer positional encodings and is therefore expected to be a partial mitigation rather than a full solution.

\subsection{Streaming Infrastructure for Species-Level Scaling}

The genus-level data-scaling question is now settled: scaling balanced reads from 50M to 250M yields only $+0.22$~pp (Section~\ref{sec:exp-scaling}), so additional \emph{genus} training data is not a productive direction within the 6-mer framework---the earlier log-linear projection into the 82--90\% range is empirically refuted (an overlapping 13-mer tokenizer, by contrast, keeps improving on the same data; Section~\ref{sec:exp-scaling}). The streaming/DDP infrastructure built for the 250M runs nonetheless remains valuable for the \emph{species}-level task (1{,}535 classes), where per-species 50M and full-scale species training are still bottlenecked by the memory-bound primary pipeline. Folding the lazy, index-based loader and multi-GPU orchestration used here into the main codebase would enable those species-level experiments without re-introducing the in-memory ceiling.

\subsection{Hierarchical Routing Improvements for Species Classification}

Per-genus 50M classifiers under oracle-genus routing reach 29.5\% species Top-1 (Section~\ref{sec:spv4-hierarchical-prelim}), exceeding the flat sp\_v4 oracle (27.4\%) and validating per-genus specialisation. Predicted routing through the 66\%--67\% genus model, however, lowers Top-1 to 15.6\%---marginally below the flat baseline (15.9\%)---because per-genus specialists do not recognise reads from genera they were not trained on.

A subgenus routing tier was investigated as a direct approach to reducing intra-classifier confusion for the five largest genera (Section~\ref{sec:subgenus-routing}). Species-level $K$-means partitioning on NT-v2 embeddings was used to define subgroups of $\sim$25--30 species. Silhouette analysis (Collinsella: 0.03) and UMAP visualisation confirm that NT-v2's 6-mer representations do not separate species within a genus, and both hard and soft subgenus routing decreased oracle Top-1 to 25.6--26.1\% (below the per-genus baseline), indicating that adding a routing tier without better representations is counterproductive.

The primary path to improving hierarchical species accuracy is therefore to improve the quality of the genus router itself. Since NT-v2's 6-mer tokenisation constrains genus accuracy to $\sim$67\%, this is best addressed through the long-$k$ pre-training direction above: a microbial foundation model with overlapping 13-mer tokenisation would simultaneously increase genus routing accuracy and provide richer intra-genus representations that could support a meaningful subgenus partition.

\subsection{Real-World Validation and Calibrated Sample-Level Estimation}

Validation on real metagenomic data from public repositories (e.g., the Human Microbiome Project~\cite{HMP2012}) is essential for assessing practical applicability. This includes (i)~robustness to sequencing errors and platform-specific quality profiles, (ii)~handling host DNA contamination, and (iii)~out-of-distribution detection for taxa absent from the training database. In parallel, the random-partition abundance evaluation in Section~\ref{sec:exp-sample} measures stability under fixed expected composition rather than generalisation across compositionally variable communities; future evaluation should therefore use sparse, asymmetric, and real microbiome-derived sample compositions, paired with (iv)~calibrated statistical tests for presence/absence above the per-sample noise floor and (v)~learned abundance correction that compensates for the observed systematic overestimation of high-abundance genera (Section~\ref{sec:sample-abundance}).

\section{Concluding Remarks}

This thesis demonstrates both the promise and the current limitations of applying general-purpose genomic foundation models to metagenomic taxonomic classification. The controlled ablation experiments yield two concrete findings that extend beyond the specific system studied here.

First, tokenization---not pre-training, depth, or further 6-mer data---sets the closed-set genus ceiling on 150~bp reads. Four architecturally different non-overlapping 6-mer models cap at 65--70\% (8-layer from scratch 65.44\%, 16-layer 67.94\%, NT-v2 LoRA 67.08\%, 29-layer from scratch 69.52\%), while overlapping 13-mer models on the same 50M reads reach 83.83\% (hashed, 16-layer), 85.63\% (exact, 16-layer), and 87.42\% (MetaTransformer, one encoder block). The previously quoted $+13.19$~pp and $+19.29$~pp ``pre-training'' gaps compared a deep pre-trained model with a one-layer from-scratch model and do not survive a depth-matched control.

Second, pre-training is not useless: at matched 16-layer depth and 6-mer tokenization it is worth $+8.94$~pp at 5M reads, and near zero (slightly negative) at 50M. That is a data-efficiency result, not a ceiling result. Combined with NT-v2's pre-training context of 2{,}048 tokens ($\sim$12{,}288~bp) versus the 25 tokens (150~bp) used here, and with a pre-training corpus and benchmark suite aimed at regulatory genomics rather than short-read taxonomy, the prescription is specific: keep pre-training as a mechanism, but pre-train on short microbial reads with overlapping long $k$-mers (or a hashed equivalent), rather than expecting a larger eukaryote-regulatory GFM to close the gap.

Realising that direction, evaluating metagenome-pre-trained BPE models that were out of scope here, and completing species-level scaling and real-community sample evaluation, represent the most promising path toward a competitive, generalisable GFM-based metagenomic classifier.
